Why is footprinting important? Hackers need to collect 
information about the target, without this phase they
will not be able to execute the attack. \todo{25-02-2026}
Thus, profiling 
of an organization. What can we get about the target 
organization publicly? 
\begin{itemize}
    \item Internet: domain names, network blocks, IP
    addresses, TCP/UDP services;
    \item Intranet: network protocols, internal domain names,
    services, access control;
    \item Remote access: remote system type, auth. mechanisms, 
    VPN;
    \item Extranet: large organizations usually have partners, 
    so their network can be connected to partner's 
    network, and hackers could get into the network 
    through the partner.
\end{itemize}

Let's start with internet footprinting. First of all, as 
a defender we need to choose the scope of our activity, we 
have to agree with the management what do we have do to, 
like profiling the entire organization or only a part. 
For the full organization, it may be worth but also very 
time consuming, as it might need a lot of employes. 
The relation between us an the management needs to be clear. 
All activities related to penetration testing and ethical 
hacking needs to be accepted by the management. \newline 
So, the second step is getting proper authorization. A 
proper funding is necessary to do this step. \newline 
\newline Let's now see the company web pages. They have public pages, 
what is shown, and private, they're not shown but they're 
still there. By looking at web pages, it's possible to 
identify security configurations, asset inventory, and 
other documents. We can also find the HTML source code, 
it gives information specifically in comments, which could 
cointain important elements like passwords. As a defender, 
we can query the web pages as much as we want, while as
attackers we'll have to use some mirroring tools, like 
\texttt{wget}, doing then an analysis offline, on a local
copy. Also, there are hidden files, and we can use 
tools like \texttt{OWASP}. After all these steps, we need 
to repeat this process for all websites of the company, like
\texttt{www1, www2, web, test}, as well as VPN sites. \newline 
After we've done this, we need to consider related 
organizations. Attacks like social engineering could 
be portraided. \newline
\newline Next step is collecting employes information. Information 
like their names, email addresses, phone numbers, other 
personal details (like home addresses), social network 
details, etc. \newline
\newline Alongside this, we can take a look at job 
postings and resumes. Also, at how the company is doing, 
like the economical state of it. Look for merges, acquisitions, 
scandals, outsourcing. For example, if we deal with merge 
and acquisition, there needs to be a merge of two IT 
systems, which could be problematic and messy. As a consequence, 
security decreases. In addition to this, we have the 
human factor, with a lot of work to do, low morale, etc. \newline
We can also take a look at archived information about the 
organization, cached and maintained in other websites. 
Services like \texttt{archive.org} does the work for us. \newline
\newline What is the counter measure for all of this? We need 
to periodically review and remove public sensitive data, as 
much as possible. \newline
\newline Once another source is given by WHOIS and DNS 
enumeration. Internet is completely distributed, there are 
independent routers everywhere. However, the management
is centralized. The ICANN organization deals with the 
management part, assigning and guarateeng correspondences 
between names and physical addresses. This database of 
logical names and IPs is also distributed, and is handled
by two services, WHOIS and ??. By querieng this services, 
we can get a lot of information about the organization, like 
IP addresses, name, mail services domains, etc. \newline
\newline Also the information contained in DNS servers 
could be quite useful. DNS could be queried (but it could 
be slow), or perform DNS zone transfer. This operation is 
useful to transfer information to increase reliability. 
However, when this takes place all the information 
is transfered between one machine to one another, and any 
attacker can do that. A specific tool is given by \texttt{dnsrecon}.
Another command is given by \texttt{nslookup}. \newline
What are the countermeasures? We can restrict transfers 
only to authorized servers, use a firewall, denying all
unauthorized connections to port TCP 53, and also not to 
put relevant information on HINFO records. \newline
\newline After we've collected all this data, we can perform 
network reconnaissance. It allows an hacker to identify 
network topology, and requires the sending of packets. 
Tools like \texttt{traceroute} can be used. The countermeasures 
are given by intrusion detection systems, blocking 
ICMP packets and UPD traffic. 

\subsubsection{Do the homework!}